% GENERATED FILE. Edit canonical lessons, then run npm run generate:books.
% canonical-source-hash: fnv1a64:bc1cd9a41d110fed
% canonical-lessons: GE-C136-wollen, GE-C136-wunsch, GE-C136-hoffen, GE-C136-bestellen, GE-C136-lieber

\chapter{To want, Wish, To hope, To order, Rather}
\label{ch:ge-to-want-wish-to-hope-to-order-rather}
\hlchaptermodality{136}

\begin{chapteropening}
\textbf{By the end of this chapter:} \emph{I can say to want, a wish, to hope, to order and rather.}

Complete the last lesson of chapter 136: I can say to want, a wish, to hope, to order and rather.
\end{chapteropening}

\section[wollen]{wollen — to want}

\label{lesson:GE-C136-wollen}



\begin{quote}
Before the new one: say the German for cool, then the German for soft.
\end{quote}

\subsection*{You'll want to know: wollen}
\textbf{wollen} — \textquotedblleft{}to want\textquotedblright{}.

\begin{cousinweb}
\textbf{wollen} goes back to Old High German \emph{wellen} $\leftarrow$ Proto-Germanic \emph{*wiljaną}, \textquotedblleft{}to want\textquotedblright{} — inherited, the same verb as English \emph{will}. (EtymDB 2.1)
\end{cousinweb}

\begin{grammarlens}[title={What you've built}]
One more word for this chapter.
\end{grammarlens}

\subsection*{Your turn}
\begin{itemize}
\raggedright
  \item \emph{Say it:} \emph{wollen}
  \item \emph{Say it:} \emph{wollen}, once more
  \item \emph{Say it:} say \emph{wollen}
  \item \emph{Recall:} say the German for cool, then the German for soft, then say \emph{wollen} again
\end{itemize}

\begin{tcolorbox}[breakable,colback=teal!4,colframe=teal!35!black,title={Before you move on}]
What does wollen mean? (\textquotedblleft{}To want\textquotedblright{}.) Say it once more.
\end{tcolorbox}
\section[der Wunsch]{der Wunsch — a wish}

\label{lesson:GE-C136-wunsch}



\begin{quote}
Before the new one: say the German for soft, then the German for to want.
\end{quote}

\subsection*{You'll want to know: der Wunsch}
\textbf{der Wunsch} — \textquotedblleft{}a wish\textquotedblright{}.

\begin{cousinweb}
\textbf{Wunsch} $\leftarrow$ Old High German \emph{wunsc} $\leftarrow$ Proto-Germanic \emph{*wunskō}, \textquotedblleft{}wish, desire\textquotedblright{} — inherited, the same word as English \emph{wish}. (EtymDB 2.1)
\end{cousinweb}

\begin{grammarlens}[title={What you've built}]
One more word for this chapter.
\end{grammarlens}

\subsection*{Your turn}
\begin{itemize}
\raggedright
  \item \emph{Say it:} \emph{der Wunsch}
  \item \emph{Say it:} \emph{der Wunsch}, once more
  \item \emph{Say it:} say \emph{der Wunsch}
  \item \emph{Recall:} say the German for soft, then the German for to want, then say \emph{der Wunsch} again
\end{itemize}

\begin{tcolorbox}[breakable,colback=teal!4,colframe=teal!35!black,title={Before you move on}]
What does der Wunsch mean? (\textquotedblleft{}A wish\textquotedblright{}.) Say it once more.
\end{tcolorbox}
\section[hoffen]{hoffen — to hope}

\label{lesson:GE-C136-hoffen}



\begin{quote}
Before the new one: say the German for to want, then the German for a wish.
\end{quote}

\subsection*{You'll want to know: hoffen}
\textbf{hoffen} — \textquotedblleft{}to hope\textquotedblright{}.

\begin{cousinweb}
\textbf{hoffen} $\leftarrow$ Middle High German \emph{hoffen}, borrowed from Low German; EtymDB 2.1 lists it beside English \emph{hope} and Dutch \emph{hopen}.
\end{cousinweb}

\begin{grammarlens}[title={What you've built}]
One more word for this chapter.
\end{grammarlens}

\subsection*{Your turn}
\begin{itemize}
\raggedright
  \item \emph{Say it:} \emph{hoffen}
  \item \emph{Say it:} \emph{hoffen}, once more
  \item \emph{Say it:} say \emph{hoffen}
  \item \emph{Recall:} say the German for to want, then the German for a wish, then say \emph{hoffen} again
\end{itemize}

\begin{tcolorbox}[breakable,colback=teal!4,colframe=teal!35!black,title={Before you move on}]
What does hoffen mean? (\textquotedblleft{}To hope\textquotedblright{}.) Say it once more.
\end{tcolorbox}
\section[bestellen]{bestellen — to order (food, a taxi)}

\label{lesson:GE-C136-bestellen}



\begin{quote}
Before the new one: say the German for a wish, then the German for to hope.
\end{quote}

\subsection*{You'll want to know: bestellen}
\textbf{bestellen} — \textquotedblleft{}to order (food, a taxi)\textquotedblright{}.

\begin{cousinweb}
\textbf{bestellen} is \emph{be-} + the verb \textquotedblleft{}to put, to place\textquotedblright{}, Old High German \emph{stellen} $\leftarrow$ Proto-Germanic \emph{*stallijaną} (EtymDB 2.1): to arrange for something to be put in place.
\end{cousinweb}

\begin{grammarlens}[title={What you've built}]
One more word for this chapter.
\end{grammarlens}

\subsection*{Your turn}
\begin{itemize}
\raggedright
  \item \emph{Say it:} \emph{bestellen}
  \item \emph{Say it:} \emph{bestellen}, once more
  \item \emph{Say it:} say \emph{bestellen}
  \item \emph{Recall:} say the German for a wish, then the German for to hope, then say \emph{bestellen} again
\end{itemize}

\begin{tcolorbox}[breakable,colback=teal!4,colframe=teal!35!black,title={Before you move on}]
What does bestellen mean? (\textquotedblleft{}To order (food, a taxi)\textquotedblright{}.) Say it once more.
\end{tcolorbox}
\section[lieber]{lieber — rather, preferably}

\label{lesson:GE-C136-lieber}



\begin{quote}
Before the new one: say the German for to hope, then the German for to order.
\end{quote}

\subsection*{You'll want to know: lieber}
\textbf{lieber} — \textquotedblleft{}rather, preferably\textquotedblright{}.

\begin{cousinweb}
\textbf{lieber} is the comparative of the word for \textquotedblleft{}dear\textquotedblright{}, Old High German \emph{liob} $\leftarrow$ Proto-Germanic \emph{*leubaz} (EtymDB 2.1): \textquotedblleft{}more dear\textquotedblright{}. English \emph{lief}, as in \textquotedblleft{}I'd as lief\textquotedblright{}, is the same word.
\end{cousinweb}

\begin{grammarlens}[title={What you've built}]
One more word for this chapter.
\end{grammarlens}

\subsection*{Your turn}
\begin{itemize}
\raggedright
  \item \emph{Say it:} \emph{lieber}
  \item \emph{Say it:} \emph{lieber}, once more
  \item \emph{Say it:} say \emph{lieber}
  \item \emph{Recall:} say the German for to hope, then the German for to order, then say \emph{lieber} again
\end{itemize}

\begin{tcolorbox}[breakable,colback=teal!4,colframe=teal!35!black,title={Before you move on}]
What does lieber mean? (\textquotedblleft{}Rather, preferably\textquotedblright{}.) Say it once more.
\end{tcolorbox}
